\section{Zero-weight lines and integral vanishing on the host frame}
\label{sec:host-van-e}

% BEGIN VHZ1
\paragraph{VHZ1. The zero-weight line on any of the supplied frames.}
\label{claim:VHZ1}
Take an algebraic closure $\Omega/\mathbb Q$, a coefficient datum
$c$ for the prime $\ell$, $N\ge1$, $e,o_1\in\mathbb N$, and a
Dirichlet character $\psi$ at level $N$ satisfying the source
$\mathrm{HostData}(\ell,e,\psi,o_1)$. Write $O=c.O$, $k=O/\mathfrak m$
and $\Delta=\mathrm{Gal}(\Omega/\mathbb Q)$. Fix any finite set
$S$ of finite places, valuation rings $A_v$ and frame $F$ satisfying
$\mathrm{IsHostFrame}(\Omega,c,e,\psi,S,A,F)$.
Thus $\pi:\Delta\twoheadrightarrow\Gamma=G_S$ is the quotient
unramified outside $S$, $D_v$ and $I_v$ are the decomposition and
inertia groups at $A_v$, and $\mathrm{loc}_v=\pi|_{D_v}$. In local
expressions write $\eta,\Psi$ also for their pullbacks. Put
\[
 T_b=\{v_\ell\}\cup S_b,\qquad
 N_v=I_v\cap\ker\eta\cap\ker\bar\Psi,\qquad H=\ker(\eta:\Gamma\to\{\pm1\}).
\]
For every $s\ge0$, take $\theta=\eta^s$ in $O^\times$.
Let $z:\Gamma\to O$ satisfy
$z(gh)=z(g)+\theta(g)z(h)$, be locally constant modulo each
$\mathfrak m^n$, and vanish on $\mathrm{loc}_v N_v$ for
$v\notin T_b$ and on $\mathrm{loc}_v I_v$ for $v\in T_b$.
Then $\mathrm{LineVan}(F,0,\theta)$ holds in its precise sense:
there exist $0\ne d\in O$ and $m\in O$ with
$d z(g)=(\theta(g)-1)m$ for every $g$.

On $H$ the action is trivial and $z$ is additive. If $t\in I_v$
has $u=\mathrm{loc}_v t\in H$ and $v\notin T_b$, set
$b=|k^\times|$. The residue field is finite, $b\ge1$, and
$t^b\in N_v$. Thus $0=z(u^b)=b z(u)$ gives $z(u)=0$ in the
characteristic-zero domain $O$. For $v\in T_b$ this vanishing
is given directly. If $gug^{-1}\in H$, then $u\in H$ and
$z(gug^{-1})=\theta(g)z(u)$ by the cocycle identity.

The class-one argument of CVI2 applies to every frame fixed above.
Indeed $\eta\pi=\eta_4$, so $\pi^{-1}H=G_{\mathbb Q(i)}$ maps
onto $H$. Any locally constant abelian character of $H$ killing all
the specified inertia conjugates pulls back to a finite-image character
of this compact Galois group. At places over $S$ the source inertias
are the intersections with the conjugates of the chosen $I_v$;
outside $S$ their images in $\Gamma$ are trivial. The pullback is
therefore everywhere unramified; $\mathbb Q(i)$ has only complex
infinite places and its Euclidean ring $\mathbb Z[i]$ has class number
one. Unramified abelian class field theory gives the trivial character.
In particular this applies to the additive homomorphism
$z|_H\bmod\mathfrak m^n$ for each $n$. Separatedness of the DVR
gives $z|_H=0$.

The function $z$ now descends to the cosets $Q=\Gamma/H$.
Set $d=|Q|\in\{1,2\}$ and $m=-\sum_{q\in Q}z(q)$. Summing
$z(gq)=z(g)+\theta(g)z(q)$ over the permutation $q\mapsto gq$
gives $d z(g)=(\theta(g)-1)m$. This produces the required
integral scalar and primitive on all of $\Gamma$.
% END VHZ1

% BEGIN VHZ2
\paragraph{VHZ2. Scalar cocycles from the reflected adjoint action.}
\label{claim:VHZ2}
For indices $a,b\in\{0,1,2\}$ write $i=(a,b)$,
$i^*=(2-b,2-a)$, $w_i=b-a$, and $(R X)_i=X_{i^*}$.
Then $R^2=1$ and $w_{i^*}=w_i$. A matrix of weight $j$ is
supported on the entries with $w_i=j$. Let $r_0=r_0(\eta,\Psi)$
be the original diagonal lift and write $L_g=\operatorname{Ad}(r_0)(g)$.
Conjugation by $\operatorname{diag}(\Psi(g),1,\Psi(g)^{-1})$
and the negative-sign differential $X\mapsto-RX$ give
\[
 (L_gX)_i=\Psi(g)^{w_i}
 \begin{cases}X_i,&\eta(g)=1,\\-X_{i^*},&\eta(g)=-1.\end{cases}
\]
Fix $j\le0$ and an integral adjoint cocycle $\gamma$ as in
$F.\mathrm{VanE}$: $\gamma(gh)=\gamma(g)+L_g\gamma(h)$,
its values have weight $j$ and its entries have locally constant
reductions at all powers of $\mathfrak m$. It vanishes on
$\mathrm{loc}_v N_v$ outside $T_b$. At $j=0$ it vanishes on
$\mathrm{loc}_v I_v$ in $T_b$. At $j<0$ its restriction to
$\mathrm{loc}_{v_\ell}(D_{v_\ell}\cap\ker\eta)$ equals
$Y-L_gY$ for one integral matrix $Y$ of weight $j$.
Here the sign kernel inside $D_{v_\ell}$ is taken with the local
pullback of $\eta$, and $g$ stands for the image in $\Gamma$.

For $s=0,1$ define $p_{i,s}(X)=X_i-(-1)^sX_{i^*}$,
$z_{i,s}(g)=p_{i,s}(\gamma(g))$, and
$\theta_{i,s}=\Psi^{w_i}\eta^s$. The two signs in the displayed
action show that $p_{i,s}(L_gX)=\theta_{i,s}(g)p_{i,s}(X)$.
Consequently each $z_{i,s}$ is a $\theta_{i,s}$-cocycle. Its
adic continuity and the stated inertia vanishings follow by
subtraction of entries. At negative weight its local scalar
primitive is $-p_{i,s}(Y)$, because the matrix coboundary
is $Y-L_gY$ and a scalar coboundary is $(\theta_{i,s}(g)-1)m$.
Both entries vanish identically when $w_i\ne j$. For $w_i=j$ these
are exactly the hypotheses of VHZ1 if $j=0$, and of VHN3 if $j<0$.
Thus on every entry and for each $s=0,1$ we have
\[
  c_{i,s}z_{i,s}(g)=(\theta_{i,s}(g)-1)m_{i,s}\qquad(g\in\Gamma)
\]
with $0\ne c_{i,s}\in O$, $m_{i,s}\in O$. Off the weight support
take $(c_{i,s},m_{i,s})=(1,0)$.
% END VHZ2

% BEGIN VHZ3
\paragraph{VHZ3. One integral matrix primitive and one denominator.}
\label{claim:VHZ3}
Let $A=\prod_{i,s}c_{i,s}\ne0$ and form the two integral matrices
$(M_s)_i=(A/c_{i,s})m_{i,s}$, $s=0,1$, using the product of all
the other factors for $A/c_{i,s}$. They have weight $j$.
Set $P_0=M_0-RM_0$ and $P_1=M_1+RM_1$. Since $z_{i^*,0}=-z_{i,0}$
and $z_{i^*,1}=z_{i,1}$, subtracting the two reflected identities
at $s=0$ and adding them at $s=1$ gives, on the weight support,
\[
 \begin{aligned}
 (\Psi(g)^j-1)(P_0)_i&=2A z_{i,0}(g),\\
 (\Psi(g)^j\eta(g)-1)(P_1)_i&=2A z_{i,1}(g).
 \end{aligned}
\]
The two sides vanish off this support. These equations include the
fixed entries $i=i^*$, where $z_{i,0}=0$ and
$z_{i,1}=2\gamma(g)_i$. As $RP_0=-P_0$ and $RP_1=P_1$, the
action in VHZ2 gives $L_gP_0=\Psi(g)^jP_0$ and
$L_gP_1=\Psi(g)^j\eta(g)P_1$. Put $X=-P_0-P_1$. Adding the two
identities entrywise and using $z_{i,0}+z_{i,1}=2\gamma_i$ yields
\[
                 (4A)\gamma(g)=X-L_gX\qquad(g\in\Gamma).
\]
Both $X$ and $4A$ belong to the original coefficient lattice and
ring, respectively, and $4A\ne0$ since $O$ is a characteristic-zero
domain. We have also treated $j<-2$, where the weight space is zero.
The denominator may depend on the frame, the weight and the cocycle.
% END VHZ3

% BEGIN VHZ4
\paragraph{VHZ4. The host vanishing installation.}
\label{claim:VHZ4}
Take the real constant $C_K$ and $\mathrm{LemmaK}(C_K)$, published
Poitou--Tate duality, local finiteness and local Tate duality in
both the $K$ and the ordinary coefficient interfaces, and local
Euler characteristic in the ordinary interface. These give VHN3
at every negative weight because $\mathrm{HostData}$ supplies
$1\le e$. The zero-weight argument VHZ1 uses the supplied frame
itself. VHZ2--3 therefore give $F.\mathrm{VanE}$ for every
$\Omega,\ell,c,N,e,\psi,o_1,S,A,F$ in the order and at the scope
of VHZ1. This is the literal $\mathrm{HostVanE}$:
for each admissible matrix cocycle there exist
$a\in O\setminus\{0\}$ and $X\in M_3(O)$ with
$a\gamma(g)=X-L_gX$ on $G_S$. It applies, in particular, to the
actual Case V frame used in CVI2.

In the outer source formulation $\mathrm{LitN}$ supplies the
normalized Poitou--Tate theorem and the specified local fields of
its cohomology record. Its normalized theorem supplies published
Poitou--Tate duality. Together with the supplied $K$-bridge those
inputs give an existential $C_K$ with $\mathrm{LemmaK}(C_K)$;
choosing that witness gives the implication
$\mathrm{LitN}\Rightarrow (K.\mathrm{KBridge}\Rightarrow
\mathrm{HostVanE})$. This establishes the selected installation
at its given cohomological inputs.
% END VHZ4
